\section{Implementation}
\label{sec:impl}

\diffkv{} is $\approx\!1.6$k lines of Python over MLX, plus the surrounding runtime and
serving code. We highlight the decisions that affect the measured behavior.

\paragraph{Attention patching.}
The wrapper rebinds \texttt{qwen2.Attention.\_\_call\_\_} to \texttt{attention\_forward} and
attaches the manager and a layer index to each attention module. This keeps the rest of the
\texttt{mlx\_lm} model---embeddings, MLP, norms, the LM head---untouched, so \diffkv{} composes
with the stock int4 model and its tokenizer without retraining or reconfiguration.

\paragraph{Why the SVD runs in NumPy.}
\texttt{mx.linalg.qr} is GPU-only in current MLX and the nonzero form of \texttt{mx.where}
does not accept the streaming kwargs the compressor needs, so the factorization is done in
NumPy on the CPU. On unified memory this is nearly free of data-movement cost: the MLX array
is materialized in place (\texttt{mx.eval}) and viewed as a NumPy array over the same bytes.
The factorization fires at most once per $256$ tokens, so its few-millisecond cost is
amortized.

\paragraph{Dtypes and precision.}
KV state, factors, and anchors are fp16; the SVD normalization scale is fp32. This halves the
store relative to fp32 defaults and matches the precision the model already runs at. The
randomized SVD itself is computed in fp32 and the factors cast to fp16 on storage.

\paragraph{Memory hygiene.}
Lazy evaluation is forced (\texttt{mx.eval}) at block boundaries to keep the graph bounded,
and \texttt{mx.clear\_cache} is called at the prefill$\to$decode transition and after block
flushes. Without these, MLX retains the peak GQA-expanded activations and the apparent
footprint balloons; with them, decode memory tracks the store.

\paragraph{Configuration and presets.}
The benchmarked configuration is \texttt{rank=16}, \texttt{block\_size=256},
\texttt{recency\_window=512} (so the dense buffer is $768$), \texttt{max\_blocks=256}, int4
weights, ``mid'' preset. Most knobs are exposed as both config keys and environment variables;
notably \texttt{DIFFKV\_COMPRESSED\_DECODE} selects the compressed kernel (default) versus
exact full-KV decode, \texttt{DIFFKV\_MAX\_BLOCKS} caps the pool, and
\texttt{DIFFKV\_ENGAGE\_THRESHOLD} sets the recency window. Three presets (low/mid/high) trade
chunk size, dense budget, and cache policy.

\ifext
\paragraph{Generation loop.}
The wrapper's generation loop supports greedy and temperature/top-$p$ sampling with a
repetition penalty and an n-gram loop detector (on a repeat ratio $\ge 0.35$ it widens the
penalty window and, if unrecovered after $40$ tokens, forces a stop). For the benchmarks we
use pure greedy decoding (\texttt{argmax}) for a fixed $128$ tokens with EOS ignored, so
throughput is directly comparable across engines.
\fi

\paragraph{The (inactive) grounding module.}
The repository includes a semantic-retrieval layer---a factual store with entity binding, a
graph query router, and logit-biasing/masking for grounded generation. In the MLX runtime as
benchmarked, the manager's \texttt{get\_srl\_state} returns \texttt{None} and the per-session
SRL map is never populated, so every factual-bias, transition-bias, and constrained-decoding
branch is skipped. We mention it because it is prominent in the code, but it contributes
nothing to the results in this report; evaluating it is future work
(\S\ref{sec:limitations}).
